\paragraph{The remaining positive-order momentum sources.}
The positive-order convection terms in the swirl and axial equations have the common form
\[
 {\cal S}^{\delta,\beta}_{ij}(U_i,F_j)
 =W_i(D_X+\delta)F_j+
 U_i[G\partial_\eta+B(\beta+\lambda_j-D_X)]F_j,\qquad
 V_i=XW_i,\quad
 W_{i,a}=\alpha_{i,a}B(U_i)_a-\frac G{a+1}\partial_\eta(U_i)_a,
\]
where $\alpha_{i,a}=1-(D+\lambda_i)/(a+1)$,
$B=2\eta/L$, $G=(1-\eta^2)/L$, $i+j=n$,
and $(\delta,\beta)=(1,b)$ in swirl or $(0,c)$ in axial
momentum. Here $b=-A-\frac12$, $c=-A$. For $a\ge i$,
$0\le\alpha_{i,a}\le1$; furthermore $|\beta|<2$.

\begin{theorem}[Uniform positive-order momentum products]
For $0<h<1/100$, $R>0$, $E=[-1,1]$ and $0<\rho\le1/4$,
\[
 \left\|\frac12J_\nu {\cal S}_{ij}^{\delta,\beta}(U_i,F_j)\right\|_{I_{n+1}}
 \le C_{\rm mom}\|U_i\|_{I_i}\|F_j\|_{I_j},
 \qquad C_{\rm mom}=\frac{64^2\,25\,R}{2}
 \left(20+\frac{16}{\rho}\right).
\]
This includes both $\nu=1$ and $\nu=2$, with a constant
independent of the two expansion orders.
\end{theorem}
\begin{proof}
The analytic coefficient algebra has multiplication constant $64$.
The radial-degree-zero functions $B,G$ have norms at most $25$:
on $|\eta|\le3$, $|B|\le300/41$ and $|G|\le500/41$;
Cauchy's inequality on radius-two discs about $[-1,1]$ and
$(m+1)^3/8^m\le1$ also bound $B',G'$ by $25$.
Use the unsigned Leibniz convolution of $U_i,F_j$.
A source with one derivative on either factor is dominated
by the total first-derivative convolution because
$\binom{s}{u}\le\binom{s+1}{u+1}$ and
$\binom{s}{u}\le\binom{s+1}{u}$.
An Euler derivative on one factor costs at most the total
radial degree $N$. The exact regular-singular inverse
multipliers satisfy
\[
 T_0\le 4R/(N+1)^2,\quad T_1\le8R/[\rho(N+1)],
 \quad NT_0\le4R,\quad (N+1)T_1\le8R/\rho.
\]
The bound $|\beta+\lambda_j-a_j|\le
(1+|\beta|)(N+1)$ holds for $a_j\ge j$.
Multiplication by $B$ or $G$ commutes with $J_\nu$ because
they do not depend on $X$, and costs a second factor $64$.
The four source types give, before the factor $1/2$,
$64^2\,25 R[4(1+\delta)+16/\rho+4(1+|\beta|)]$,
which is bounded as stated.
\end{proof}

\paragraph{Complete two-shift delayed momentum forcing.}
For $\beta\in\{b,c\}$, $n\ge2$, $a\ge n-1$, put
$r=a+1$, $\theta=\beta+\lambda_{n-1}-a$, $\theta_2=\theta-D$.
The \emph{full} variable-coefficient normal form is
\[
 Z_{\beta,n-1}^{[2]}F_a=
 (\theta GB'+\theta\theta_2B^2)F_a+
 [GG'+(\theta+\theta_2)GB]F'_a+G^2F''_a.
\]
Since $|\theta|,|\theta_2|\le3r$, its three
radial-normalized coefficient norms are bounded by
$12(64)25^2$, $7(64)25^2$ and $64\,25^2$.
Using $r^2T_0\le4R$, $rT_1\le8R/\rho$ and
$T_2\le8R/\rho^2$ gives, for $\nu=1,2$,
\begin{equation}\label{eq:direct-mom-uniform}
 \left\|\frac12J_\nu Z_{\beta,n-1}^{[2]}F\right\|_{I_n}
 \le\frac{64^2\,25^2R}{2}
 \left(48+\frac{56}{\rho}+\frac8{\rho^2}\right)
 \|F\|_{I_{n-1}}.
\end{equation}
This proves an $n$-independent bound for the complete
\emph{direct} delayed swirl and axial momentum terms.

\paragraph{A pressure-norm obstruction to the stated majorant.}
The preceding theorems bound the velocity forcing. They do
not justify the unchanged normalized amplitude in equation
(52), which includes $\|\Pi_n\|_{I_{n-1}}/(n+1)$.
The isolated delayed-viscosity pressure primitive is not
a bounded map on that same weighted space.

\begin{proposition}[Explicit pressure-primitive counterexample]
Take $h=1/200$, $\rho=1/4$ and $R>0$.
With input $F=X/(1+4i\eta)\in I_1$,
the isolated delayed pressure source produces
$\Pi_3^{\rm del}=X^2P(\eta)$ where
$\|F\|\le8R$ but $\|\Pi_3^{\rm del}\|_{I_2}=\infty$.
\end{proposition}
\begin{proof}
Put $r=2$, $\lambda=2h$,
$W=B(r-D-\lambda)f/r-Gf'/r$,
$f=(1+4i\eta)^{-1}$, and
$Z_q^{(r)}=G\partial_\eta+B(q-r)$.
The exact pressure primitive from (24), (28), (27) is
\[
 P=\frac1{2r}Z_{\lambda-D}^{(r)}Z_\lambda^{(r)}W.
\]
At $\zeta=i/4$ the coefficient of its fourth-order pole is
\[
 \lim_{\eta\to\zeta}(\eta-\zeta)^4P(\eta)
 =-921187500i/4103684801\ne0.
\]
All other coefficient singularities are farther from the
origin. Thus $|P^{(m)}(0)|\asymp4^m m!\,m^3$.
But $w_{2,m}\asymp R^{-2}4^m m!\,m^2$, whereas
$w_{1,m}=R^{-1}4^m m!(m+2)/[8(m+1)]$.
In particular the precise scalar gap is
\[
 \frac{w_{1,m+3}}{4w_{2,m}}
 =\frac{27R}{4\rho^3}(m+5)
       \left(\frac{m+1}{m+4}\right)^2\longrightarrow\infty.
\]
This establishes the claimed unboundedness.
\end{proof}
This operator-level counterexample does not rule out the
specific recursively generated solution. A possible repair is
to eliminate the bare pressure norm from (52), use the
regularised pressure-to-velocity composite throughout the
induction, and reconstruct $\Pi_n$ in a weaker pressure
space, for example using $w_{a,m+1}$ in place of $w_{a,m}$.
The full recurrence (53) and Theorem~1.1 remain unproved
until that repair and all physical-coordinate reconstructions
are independently verified.